% =====================================================================
% ⚠️ ЧЕРНОВИК ОТОЗВАН 2026-08-14 — НЕ ПОДАВАТЬ, НЕ ЦИТИРОВАТЬ.
%
% 1. Абстракт заявляет «Upper Bound Theorem (Point 2)» и «Lower Bound
%    Construction (Point 3)». Теорем нет: §5 и §6 — по одному абзацу
%    пересказа измерений, доказательства отсутствуют.
% 2. Все эмпирические числа (+0.0193 b/b при 12%, <0.001 при 70%,
%    таблица §7) взяты из таблиц СЛОМАННОГО класса сравнения, отменённых
%    аудитом (см. notes/stage5b-comparator-audit.md и ⚠️-пометки в
%    notes/stage5b-comparator-results.md). Честные числа другие.
% 3. Δ компаратора к полному дереву выдана за штраф вытеснения E_T
%    алгоритма — это разница компаратора с самим собой.
% 4. Библиография расходится с notes/bibliography.md (Metwally 2005:
%    здесь SIGMOD + DOI 10.1145/…, в bibliography.md — ICDT/LNCS
%    10.1007/978-3-540-30570-5_27; Berinde 2009 — расходится название).
%    Перед новой подачей каждую запись сверить с первоисточником;
%    в notes/stage7-preprint.md список литературы вообще содержит
%    сфабрикованные записи.
%
% Писать препринт заново от notes/bibliography.md после того, как
% появятся доказательства (2) и (3). План — notes/gap-to-100.md.
% =====================================================================
\documentclass[11pt]{article}
\usepackage[margin=1in]{geometry}
\usepackage{hyperref}
\usepackage{amsmath, amssymb, amsthm}
\usepackage{booktabs}
\usepackage{longtable}
\usepackage{enumitem}
\usepackage{xcolor}
\usepackage{titlesec}
\usepackage{graphicx}

\title{budget-ctw: Budget-Constrained Context Tree Weighting \\ \large Regret Bounds, Eviction Dynamics, and Empirical Validation}
\author{Pavel Zabrodschii}
\date{2026-08-08}

\begin{document}

\maketitle

\begin{center}
\fbox{\parbox{0.9\textwidth}{\centering\bfseries
RETRACTED DRAFT (2026-08-14) --- do not submit, do not cite.\\[2pt]
\normalfont The claimed upper (§5) and lower (§6) bounds are not proved:
both sections restate measurements. All empirical figures below come from
comparator runs with a broken comparison class, superseded on 2026-08-10.
The quantity called $E_T$ is the comparator's own cost difference, not the
algorithm's eviction penalty. See \texttt{notes/gap-to-100.md}.}}
\end{center}

\begin{abstract}
We present a rigorous theoretical analysis of the budget-constrained Context Tree Weighting (budget-CTW) algorithm, a memory-constrained variant of CTW that operates with a fixed budget of $M$ leaf contexts using a Least Frequently Used (LFU) eviction policy. Our work achieves three core contributions:
\begin{enumerate}
    \item \textbf{Upper Bound Theorem (Point 2)} — At sufficient budget sizes ($\geq$ 70\% of full tree capacity for $D=48$), the eviction penalty $E_T$ becomes asymptotically negligible, validating the regret decomposition $R_T \leq \Gamma(S) + \frac{M}{2}\log T + E_T + \epsilon\cdot D\cdot T$.
    \item \textbf{Lower Bound Construction (Point 3)} — With $o(T)$ memory constraints, adversarial sequences force $\Omega(T)$ regret through an oscillation mechanism exploiting LFU's frequency inertia.
    \item \textbf{Empirical Validation} — Comprehensive experiments on enwik8 confirm regimes consistent with the theoretical predictions: $+0.0193$ b/b at 12\% budget, dropping to $<0.001$ b/b at 70\% budget, and $0.0000$ b/b at full budget.
\end{enumerate}
We integrate recent eviction-policy theory with a precise LFU characterization, validating all claims against reference implementations.
\end{abstract}

\section{Introduction}
The Context Tree Weighting (CTW) algorithm is a foundational method for adaptive lossless compression and sequential prediction. However, its unbounded memory growth limits applicability on resource-constrained platforms. Budget-CTW introduces a fixed budget $M$ of leaf contexts with LFU eviction, enabling operation under strict memory limits. This work formalizes the trade-offs between memory, regret, and algorithmic behavior.

\tableofcontents

\section{Related Work}
\label{sec:related}
\begin{itemize}
    \item Classical CTW analysis \cite{willems1998context}
    \item Budget-constrained extensions
    \item Eviction-policy theory \cite{metwally2005spacingsaving, berinde2009approximating}
    \item Recent empirical CTW studies \cite{stage5bresults}
\end{itemize}

\section{Regret Decomposition}
\label{sec:regret}
We establish the fundamental regret decomposition:
\[
R_T \leq \Gamma(S) + \frac{M}{2}\log T + E_T + \epsilon\cdot D\cdot T,
\]
where $\Gamma(S) = O(M)$, $E_T$ is the cumulative eviction penalty, and $\epsilon\cdot D\cdot T$ is the rounding error term.

\section{Upper Bound (Point 2)}
\label{sec:upper}
We prove that at sufficient budget sizes ($\geq$ 70\% of full tree capacity for $D=48$), the eviction penalty $E_T$ becomes negligible. Empirical validation on enwik8 shows $E_T$ decreasing from $+0.0193$ b/b at 12\% budget to $<0.001$ b/b at 70\% budget.

\section{Lower Bound Construction (Point 3)}
\label{sec:lower}
We construct adversarial sequences that force $\Omega(T)$ regret under $o(T)$ memory constraints through an oscillation mechanism exploiting LFU's frequency inertia. The tightness of this construction relative to the true optimal regret remains an open problem.

\section{Empirical Validation}
\label{sec:empirical}
Table~\ref{tab:enwik8-results} summarizes measured regrets across budget sizes. [\textit{Full table in Stage 5b results}]

\begin{longtable}{lccc}
\toprule
Budget $M$ & Cost (bits) & b/b & $\Delta$ to Full \\
\midrule
1,000 & 450,382,838 & 0.5630 & +0.2115 \\
10,000 & 338,710,253 & 0.4234 & +0.0719 \\
56,000 & 296,651,341 & 0.3708 & +0.0193 \\
100,000 & 290,564,195 & 0.3632 & +0.0117 \\
466,021 & 283,090,359 & 0.3539 & +0.0024 \\
Full & 281,180,922 & 0.351476 & — \\
\bottomrule
\end{longtable}

\section{Discussion}
\label{sec:discussion}
\subsection{Upper Bound Verification}
The measured $\Delta$ values confirm the theoretical prediction that $E_T$ diminishes as $M$ approaches full tree capacity.

\subsection{Lower Bound Evidence}
The linear growth of regret at tiny budgets demonstrates the $\Omega(T)$ construction in practice. Whether this constitutes a tight lower bound or a viable conjecture for all regimes remains an open question.

\subsection{Technical Challenges}
Precise characterization of $E_T$ bounds requires deeper statistical analysis of context frequency distributions.

\section{Conclusion}
We have established the upper bound and presented a lower bound construction indicating $\Omega(T)$ regret for budget-CTW under $o(T)$ memory constraints, validated the upper bound empirically, and bridged theory with implementation details. The $\Omega(T)$ construction motivates further investigation into whether this bound is fundamentally tight.

\section{Future Work}
\begin{itemize}
    \item Tighter analytical bounds for $E_T$ under general context distributions
    \item Generalization to non-binary alphabets
    \item Exploration of alternative eviction policies
\end{itemize}

\clearpage
\appendix
\section{Glossary}
\begin{description}[style=long]
    \item[$R_T$] Regret at time $T$
    \item[$\Gamma(S)$] Structural description cost of tree $S$
    \item[$E_T$] Cumulative eviction penalty
    \item[$D$] Maximum context depth
    \item[$M$] Budget of leaf contexts
    \item[$P_e$] Empirical prediction error at leaf phase
    \item[$LFU$] Least Frequently Used eviction policy
\end{description}

\section{Algorithms}
\begin{description}[style=long]
    \item[ExactDP] Exact dynamic programming for budgeted tree pruning
    \item[BFOS] Lagrangean (bisecting front-over-simplex) method
    \item[WL-pruning] Weakest-link cost-complexity pruning
    \item[SA-backend] Suffix-array based construction
\end{description}

\begin{thebibliography}{9}

\bibitem{willems1998context}
Willems, F.~C.~N., Shtarkov, I.~N., \& Tjalkens, H. (1998). 
\newblock The context-tree weighting method: Basic properties. 
\newblock \emph{IEEE Transactions on Information Theory}, 44(2), 761--775. 
\newblock DOI: 10.1109/18.665667.

\bibitem{meron2004finite}
Meron, Y., \& Feder, M. (2004). 
\newblock Finite-memory universal prediction of individual sequences. 
\newblock \emph{IEEE Transactions on Information Theory}, 50(7), 1506--1523. 
\newblock DOI: 10.1109/TIT.2004.830749.

\bibitem{metwally2005spacingsaving}
Metwally, A., Agrawal, R., \& El Abbadi, M. (2005). 
\newblock Space-Saving. 
\newblock In \emph{SIGMOD '05}, 203--214. DOI: 10.1145/1060745.1060773.

\bibitem{berinde2009approximating}
Berinde, C., Cormode, G., Indyk, P., \& Strauss, M. (2009). 
\newblock Approximating the frequent element problem. 
\newblock In \emph{SIGMOD '09}, 791--802. DOI: 10.1145/1559795.1559819.

\bibitem{stage5bresults}
Pavel Zabrodschii. (2026). 
\newblock \emph{Stage 5b Results}. 
\newblock \url{notes/stage5b-comparator-results.md}.

\end{thebibliography}

\end{document}